\section{Chapter~\ref{sec:dishbrain}. DishBrain}

The bench design and the game results come from one experimental study and its follow-up; the formulas for noise and for drift toward good configurations are derived in the chapter and checked with the book's model, while the learning mechanism in the dish has not been established.

{\footnotesize
\setlength{\tabcolsep}{3pt}\renewcommand{\arraystretch}{1.15}
\begin{longtable}{>{\raggedright}p{0.5cm}>{\raggedright}p{4.0cm}>{\raggedright}p{5.6cm}>{\raggedright}p{2.3cm}>{\raggedright\arraybackslash}p{1.7cm}}
\toprule
No. & Claim & Experiment and what was measured & Source & Status \\
\midrule\endfirsthead
\toprule
No. & Claim & Experiment and what was measured & Source & Status \\
\midrule\endhead
1 & Bench: about $800\,000$ mouse cortical cells or about $10^6$ human cells per chip; $26\,000$ electrodes on $8$~mm$^2$, recording up to $1024$, stimulation through $8$ & Methods of the study: MaxOne chip, $26\,000$ platinum electrodes on $8$~mm$^2$, recording of $1024$ channels at $20\,000$ samples per second; up to $32$ electrodes can technically be stimulated, $8$ could be stimulated independently; mouse cultures were used from about $10$~days & \cite{Kagan2022} & measured \\
2 & A loop with a $10$~ms clock: spikes from the motor regions move the paddle (fig.~\ref{fig:dishloop}) & Spikes are counted in $10$~ms windows; the delay from a spike to the responding stimulus is about $5$~ms, set by the hardware buffer & \cite{Kagan2022} & measured \\
3 & Input: place code (which of $8$ electrodes) and rate code, $4$--$40$~pulses/s & In the main experiment the stimulation rate rises linearly from $4$~Hz when the ball is at the far wall to $40$~Hz at the paddle; the amplitude of the ball pulses is $75$~mV; the pulses are biphasic and rectangular & \cite{Kagan2022} & measured \\
4 & Output $\Delta p=c\,(n_\uparrow-n_\downarrow)$~\eqref{eq:dishreadout}, count noise per window $1/\sqrt{RT}$ & The two-region difference rule is described in the study (with region weights by activity); the exact formula is not given. The noise estimate follows from Poisson counting and is derived in the chapter & \cite{Kagan2022}; derived in the chapter & theory \\
5 & Teacher: after a hit, a predictable stimulus of $75$~mV, $100$~pulses/s, $0.1$~s; after a miss, an unpredictable one, $150$~mV, $5$~pulses/s, $4$~s at random places and times, then a $4$~s pause & Methods: after a hit, $75$~mV, $100$~Hz, $100$~ms to all $8$ electrodes at once; after a miss, $150$~mV, $5$~Hz to random electrodes at random times for $4$~s, then $4$~s without stimulation. The culture contains no dopamine & \cite{Kagan2022} & measured \\
6 & The network acts so as to make its input predictable & The free-energy principle is a theoretical proposal from which the authors derived the protocol; the experiment agrees with it but does not distinguish it from simpler mechanisms (rows~7--8) & \cite{Friston2010,Kagan2022} & hypothesis \\
7 & If failure shakes the network and success does not, the time in a configuration is $\rho\propto1/P_{\text{miss}}$~\eqref{eq:dishdensity} (proposition~\ref{prop:dishscramble}) & Follows from flux balance in a Markov chain with symmetric kicks; derived in the chapter. There is no direct test on a culture & derived in the chapter & theory \\
8 & The ``scramble on a miss'' model learns: fraction of hits $0.21\to0.38$; with a kick after every rally $\approx0.12$, with no kicks $0.19$ (fig.~\ref{fig:dishmodel}) & Computation, script \texttt{models/dishbrain\_model.py}, $40$ model cultures with $2000$ rallies each: $0.21\to0.38$; $0.13\to0.11$; $0.19\to0.19$ & --- & model \\
9 & Runs of hits grow longer only in living cultures in the full loop; the controls do not learn & $399$ sessions of $20$~min: $9$ mouse cultures ($101$ sessions), $11$ human ($138$), $6$ medium-only chips ($80$), $20$ cultures at rest ($42$), $3$ simulations ($38$). The mean rally length over the last $15$~min exceeds that over the first $5$ only in living cultures; non-neuronal cells (HEK293T) and medium-only chips show no effect & \cite{Kagan2022} & measured \\
10 & Significant growth within a session only with full feedback; with silence in place of the stimulus, play is better by the end of the session than without feedback, but the effect is smaller & $486$ sessions over $3$ days: ``stimulus'' ($n=27$), ``silence'' ($17$), ``no feedback'' ($15$). Growth over time is significant only in the ``stimulus'' condition; at the end of the session ``silence'' outperformed ``no feedback'' with a smaller effect & \cite{Kagan2022} & measured \\
11 & The effect is small; whether the network learns for the long term is an open question & Within a session the improvement is reliable, but, as the authors themselves say, learning across sessions over several days was not robustly observed & \cite{Kagan2022} & measured \\
12 & During play the network approaches a critical regime, and the closer it gets, the better the play & Analysis of activity avalanches in the same cultures: structured input brings the network closer to the critical regime, and play quality correlates with proximity to it; but criticality alone, without information about the consequences of actions, is not enough for learning & \cite{Habibollahi2023} & measured \\
13 & The culture's ``sentience'' has not been shown & A response article by thirty researchers: behavior change in a closed loop proves neither intelligence nor sentience; no experiment shows it & \cite{Balci2023} & hypothesis \\
14 & Proposed fourth control: a predictable stimulus after a miss, of the same duration and energy (section~\ref{sec:dishspec}) & No such experiment has been done; this is the book's proposal & --- & hypothesis \\
\bottomrule
\end{longtable}}
